%%=============================================================================
%% Inleiding
%%=============================================================================

\chapter{\IfLanguageName{dutch}{Inleiding}{Introduction}}%
\label{ch:inleiding}

De Belgische zakelijke markt staat voor een ingrijpende digitale transformatie. Vanaf 1 januari 2026 zijn Belgische ondernemingen verplicht om gestructureerde B2B-e-facturen systeem-tot-systeem uit te wisselen via het Peppol-netwerk \autocite{Deloitte_2024}. Deze verplichting past binnen een bredere Europese digitaliseringsstrategie rond e-invoicing en interoperabiliteit \autocite{EuropeanCommission_2025}. Voor veel kleine ondernemingen, die vandaag nog sterk steunen op manuele administratie in Excel, e-mail en afzonderlijke softwaretools, vormt deze deadline niet enkel een compliance-uitdaging, maar ook een aanleiding om kernprocessen fundamenteel te herdenken.

Enterprise Resource Planning (ERP)-systemen worden in de literatuur omschreven als geïntegreerde informatiesystemen die bedrijfsprocessen over functionele silo's heen verbinden \textcite{Davenport1998}. Een succesvolle ERP-implementatie vereist echter meer dan louter software-installatie. Ze impliceert procesaanpassing, dataconsistentie, gebruikersacceptatie en organisatorische verandering \autocite{Nah2001}. Binnen digitale transformatie bij KMO's wordt ERP vaak beschouwd als een katalysator voor standaardisatie, schaalbaarheid en betere procescontrole \autocite{Matt2015,Vial2019}.

Deze bachelorproef onderzoekt deze transitie aan de hand van een concrete casus: Hanuman BV, een Belgische KMO gespecialiseerd in de verkoop, reparatie en het onderhoud van professionele koffieapparatuur. Door het hybride karakter van de onderneming, met zowel productverkoop als serviceactiviteiten, is er een duidelijke nood aan consistente stamdata, end-to-end procesopvolging, voorraadactualiteit en betrouwbare facturatie. Vanuit deze context wordt onderzocht hoe een gefaseerde migratie naar het open-source ERP-systeem Odoo kan bijdragen aan procesverbetering en aan een correcte voorbereiding op B2B-e-invoicing via Peppol.

\section{\IfLanguageName{dutch}{Probleemstelling}{Problem Statement}}%
\label{sec:probleemstelling}

Dit onderzoek richt zich op Belgische, servicegerichte KMO's met een beperkte personeelsomvang die hun administratieve kernprocessen willen digitaliseren, maar beperkt zijn in tijd, expertise en middelen. De primaire casus in deze bachelorproef is Hanuman BV, waarbij de zaakvoerder en administratieve medewerkers de belangrijkste stakeholders vormen.

Binnen Hanuman BV veroorzaken gefragmenteerde tools en Excel-gestuurde opvolging vertragingen en foutgevoeligheid. Typische knelpunten zijn dubbele invoer tussen verkoop, voorraad en facturatie, beperkte traceerbaarheid van herstellingen en onderhoud, manuele correcties bij facturen en betalingen, en een gebrek aan centraal inzicht in voorraad- en klantgegevens. Daarnaast moet Hanuman BV zich tijdig voorbereiden op de verplichting rond gestructureerde B2B-e-facturatie via Peppol, waardoor een werkwijze met PDF-facturen via e-mail op termijn onvoldoende wordt.

Er ontbreekt binnen deze context een concreet, stapsgewijs migratiekader waarmee een KMO de overstap naar een geïntegreerd ERP-systeem gecontroleerd kan uitvoeren. Zonder gefaseerde aanpak ontstaan risico's zoals datakwaliteitsproblemen, weerstand bij gebruikers, foutieve factuurgegevens, onduidelijke procesverantwoordelijkheden en scope creep. De uitdaging bestaat er dus niet alleen in om Odoo technisch te configureren, maar vooral om de migratie beheersbaar, reproduceerbaar en afgestemd op de bedrijfsprocessen te maken.

\section{\IfLanguageName{dutch}{Onderzoeksvraag}{Research question}}%
\label{sec:onderzoeksvraag}

De centrale onderzoeksvraag van deze bachelorproef luidt als volgt:

\begin{quote}
    \textit{Hoe kan een stapsgewijze aanpak voor de migratie naar open-source ERP (Odoo) KMO's zoals Hanuman BV helpen om doorlooptijd, foutgevoeligheid en dubbele gegevensinvoer te verminderen, en tegelijk de voorbereiding op B2B-e-invoicing via Peppol te ondersteunen?}
\end{quote}

Om deze onderzoeksvraag te beantwoorden, worden de volgende deelvragen geformuleerd:

\begin{itemize}
    \item Welke knelpunten in de AS-IS-processen, zoals verkoop, voorraadbeheer en service/herstellingen, veroorzaken bij Hanuman BV de grootste vertragingen en foutbronnen?
    \item Welke baseline-KPI's, zoals het aantal manuele stappen, doorlooptijd, correctiemomenten en openstaande posten, kunnen worden vastgesteld voor de geselecteerde processen?
    \item In welke mate kan een gefaseerde implementatie van Odoo-modules, zoals Sales, Inventory, Accounting/Invoicing, Purchase, Repair en Maintenance, deze processen geïntegreerd ondersteunen?
    \item Welke configuratie- en datavereisten, zoals stamdata, btw-gegevens, klantgegevens, facturatie-instellingen en Peppol-identificatoren, zijn cruciaal om Peppol-readiness te bereiken?
    \item Welke beperkingen blijven bestaan wanneer de voorgestelde Odoo-aanpak enkel gevalideerd wordt binnen een gecontroleerde Proof of Concept-omgeving?
\end{itemize}

\section{\IfLanguageName{dutch}{Onderzoeksdoelstelling}{Research objective}}%
\label{sec:onderzoeksdoelstelling}

Het doel van deze bachelorproef is het ontwerpen en evalueren van een praktisch, herbruikbaar migratiekader voor de overstap van een gefragmenteerde Excel-/toolomgeving naar Odoo binnen een servicegerichte Belgische KMO. Dit kader wordt gevalideerd via een Proof of Concept in Odoo en een set testscenario's gebaseerd op realistische transacties, zoals offerte--order--factuur, aankoop en ontvangst, voorraadbewegingen, repair orders, onderhoudsplanning en e-invoicingvoorbereiding.

De Proof of Concept heeft als doel om na te gaan of de ontworpen TO-BE-processen functioneel uitvoerbaar zijn binnen een gecontroleerde Odoo-testomgeving. Daarbij wordt niet gestreefd naar een volledige productie-implementatie, maar naar een onderbouwde validatie van de belangrijkste procesketens en configuratievereisten.

De succescriteria van dit onderzoek zijn:

\begin{itemize}
    \item een aantoonbare vermindering van manuele stappen en dubbele gegevensinvoer binnen de geselecteerde processen;
    \item een betere traceerbaarheid tussen verkooporders, leveringen, facturen, repair orders, serienummers en onderhoudsopdrachten;
    \item automatische voorraadactualisatie bij aankoop, verkoop en onderdelenverbruik;
    \item het kunnen voorbereiden van facturen voor gestructureerde elektronische facturatie, inclusief Peppol-demo en UBL BIS 3 XML-generatie binnen de testomgeving;
    \item het formuleren van een gefaseerd migratiekader dat bruikbaar is voor Hanuman BV en vergelijkbare servicegerichte KMO's.
\end{itemize}

De resultaten worden geïnterpreteerd binnen de grenzen van de gebruikte testomgeving. Productie-Peppol, live bankkoppelingen, historische datamigratie en langdurige gebruikersacceptatie vallen buiten de scope van deze bachelorproef.

\section{\IfLanguageName{dutch}{Opzet van deze bachelorproef}{Structure of this bachelor thesis}}%
\label{sec:opzet-bachelorproef}

In Hoofdstuk~\ref{ch:stand-van-zaken} wordt de stand van zaken besproken op basis van een literatuurstudie over ERP-implementaties in KMO-context, Odoo als modulair ERP-platform en de Belgische en Europese context rond e-invoicing en Peppol.

Hoofdstuk~\ref{ch:methodologie} licht de onderzoeksaanpak toe. Daarbij worden de casestudybenadering, de AS-IS-analyse, het ontwerp van het gefaseerde migratiekader, de configuratie van de Odoo Proof of Concept en de methode voor KPI-evaluatie besproken.

In Hoofdstuk~\ref{ch:implementatie-analyse} wordt de implementatie en analyse van de casus uitgewerkt. Daarbij worden de AS-IS-situatie, de belangrijkste knelpunten, requirements, scope, het TO-BE-ontwerp en het gefaseerde migratiekader besproken als voorbereiding op de Proof of Concept.

In Hoofdstuk~\ref{ch:poc} wordt de Proof of Concept uitgewerkt. Hierbij worden de Odoo-testomgeving, de gebruikte stamdata, de uitgevoerde testscenario's, screenshots en validatieresultaten besproken.

In Hoofdstuk~\ref{ch:evaluatie} worden de resultaten geëvalueerd aan de hand van KPI's, waaronder manuele stappen, doorlooptijd, traceerbaarheid, voorraadactualiteit en Peppol-readiness.

In Hoofdstuk~\ref{ch:conclusie} worden de conclusies geformuleerd, de centrale onderzoeksvraag beantwoord en aanbevelingen gegeven voor verdere uitrol en toekomstig onderzoek.